% ============================================================
% Title supplied by appendix wrapper.
\label{sec:advanced-operators-ii}
% ============================================================

These operators extend the earlier toolbox (horizon flow, spectral selection, traces, polar decomposition) by adding symmetry projectors, capacity/entropy functionals, graded topological invariants, and stability indices.

% ------------------------------------------------------------
\subsection{Representation Projectors (Isotypic Decomposition under Horizon Symmetry)}
\label{subsec:representation-projector}
% ------------------------------------------------------------

When the observable subspace carries a symmetry, representation projectors decompose it into invariant components.

\begin{definition}[Group action on observable subspace]
\label{def:group-action}\lean{CoherenceCalculus.AppendixOperators.GroupAction}
Let $G$ be a finite group acting unitarily on the observable subspace $\cH_{\mathrm{obs}} := \CCPi{h} \cH$ via a representation $U: G \to U(\cH_{\mathrm{obs}})$, meaning $U(g)$ is unitary for each $g \in G$ and $U(gh) = U(g)U(h)$.
\end{definition}

\begin{definition}[Isotypic projector]
\label{def:isotypic-projector}\lean{CoherenceCalculus.AppendixOperators.AppendixIsotypicProjector}
For an irreducible representation $\rho$ of $G$ with character $\chi_\rho$ and dimension $d_\rho$, define the \emph{isotypic projector}
\[
P_\rho := \frac{d_\rho}{|G|} \sum_{g \in G} \overline{\chi_\rho(g)}\, U(g).
\]
(For unitary representations, $\overline{\chi_\rho(g)} = \chi_\rho(g^{-1})$.)
This projects onto the $\rho$-isotypic component of $\cH_{\mathrm{obs}}$.
\end{definition}

\begin{proposition}[Properties of isotypic projectors]
\label{prop:isotypic-properties}\lean{CoherenceCalculus.AppendixOperators.IsotypicProperties}
The isotypic projectors satisfy:
\begin{enumerate}[label=(\roman*)]
\item $P_\rho^2 = P_\rho = P_\rho^*$ (orthogonal projector).
\item $P_\rho P_{\rho'} = 0$ for distinct irreps $\rho \ne \rho'$.
\item $\sum_{\rho} P_\rho = I$ on $\cH_{\mathrm{obs}}$ (sum over irreps appearing in the representation).
\end{enumerate}
\end{proposition}

\begin{proof}
These are the standard central idempotents attached to the finite-dimensional
unitary representation \(U\).

For self-adjointness, use \(U(g)^* = U(g^{-1})\) together with
\(\overline{\chi_\rho(g)} = \chi_\rho(g^{-1})\):
\[
P_\rho^*
=
\frac{d_\rho}{|G|}\sum_{g\in G} \chi_\rho(g^{-1})\,U(g)^*
=
\frac{d_\rho}{|G|}\sum_{g\in G} \chi_\rho(g^{-1})\,U(g^{-1}).
\]
Reindexing \(g \mapsto g^{-1}\) gives \(P_\rho^* = P_\rho\).

For the product, write
\[
P_\rho P_{\rho'}
=
\frac{d_\rho d_{\rho'}}{|G|^2}
\sum_{g,h\in G}
\overline{\chi_\rho(g)}\,
\overline{\chi_{\rho'}(h)}\,
U(gh).
\]
Reindex with \(k = gh\) and apply the Schur orthogonality relations for
characters. This yields
\[
P_\rho P_{\rho'} = \delta_{\rho\rho'}\,P_\rho,
\]
so in particular \(P_\rho^2 = P_\rho\) and \(P_\rho P_{\rho'} = 0\) for
\(\rho \neq \rho'\).

Finally, the finite-dimensional unitary representation decomposes as the direct
sum of its isotypic components, and the corresponding central idempotents sum
to the identity on \(\cH_{\mathrm{obs}}\). Hence
\[
\sum_\rho P_\rho = I
\]
on the observable subspace.
\end{proof}

\begin{corollary}[Invariant subspace dimension]
\label{cor:invariant-dimension}\lean{CoherenceCalculus.AppendixOperators.invariantDimension}
The dimension of the $G$-invariant subspace is
\[
\dim(\cH_{\mathrm{obs}}^G) = \rank(P_{\mathrm{triv}}) = \frac{1}{|G|} \sum_{g \in G} \Tr(U(g)),
\]
where $P_{\mathrm{triv}} := \frac{1}{|G|} \sum_{g \in G} U(g)$ is the projector onto the trivial isotypic component.
\end{corollary}

\begin{example}[Symmetric group $S_2$ on $\bbR^2$]
\lean{CoherenceCalculus.AppendixOperators.s2_projectors_example}
\label{ex:s2-projectors}
Let $G = S_2 = \{e, \sigma\}$ act on $\bbR^2$ by swapping coordinates: $U(\sigma)(x_1, x_2) = (x_2, x_1)$.
In the standard basis, $U(\sigma) = \begin{psmallmatrix} 0 & 1 \\ 1 & 0 \end{psmallmatrix}$.
The two irreps are trivial ($\chi_{\mathrm{triv}}(e) = \chi_{\mathrm{triv}}(\sigma) = 1$) and sign ($\chi_{\mathrm{sign}}(e) = 1$, $\chi_{\mathrm{sign}}(\sigma) = -1$).
\begin{align*}
P_{\mathrm{triv}} &= \tfrac{1}{2}(I + U(\sigma)) = \tfrac{1}{2}\begin{psmallmatrix} 1 & 1 \\ 1 & 1 \end{psmallmatrix}, \\
P_{\mathrm{sign}} &= \tfrac{1}{2}(I - U(\sigma)) = \tfrac{1}{2}\begin{psmallmatrix} 1 & -1 \\ -1 & 1 \end{psmallmatrix}.
\end{align*}
These project onto $\Span\{(1,1)\}$ and $\Span\{(1,-1)\}$ respectively.
\end{example}

\begin{quote}\small
\textbf{Intuition.}
Constitutive laws and admissible forms often must be invariant under symmetries.
Isotypic projectors compute the allowed (invariant or covariant) components directly.
This connects to the completion-and-selection algebra
(Section~\ref{sec:constraint-compiler}): symmetry is one of the constraints
that reduces the rule space.
\end{quote}

% ------------------------------------------------------------
\subsection{Rank-Log Functionals (Dimension Measures for Horizons)}
\label{subsec:horizon-entropy}
% ------------------------------------------------------------

Entropy functionals measure the ``representational capacity'' of a horizon.

\begin{definition}[Horizon density operator]
\label{def:horizon-density}\lean{CoherenceCalculus.AppendixOperators.horizonDensity}
For finite-dimensional $\cH$ (or when $\CCPi{h}$ has finite rank), define the normalized \emph{horizon density operator}
\[
\rho_h := \frac{\CCPi{h}}{\Tr(\CCPi{h})}.
\]
This is a density operator (positive semidefinite with trace 1) supported on $\cH_{\mathrm{obs}}$.
\end{definition}

\begin{definition}[Horizon entropy]
\label{def:horizon-entropy}\lean{CoherenceCalculus.AppendixOperators.horizonEntropy}
The \emph{horizon entropy} is the von Neumann entropy of the horizon density:
\[
S(\CCPi{h}) := -\Tr(\rho_h \log \rho_h).
\]
\end{definition}

\begin{lemma}[Entropy of a projection]
\label{lem:projection-entropy}\lean{CoherenceCalculus.AppendixOperators.projection_entropy}
If $\CCPi{h}$ is an orthogonal projection of rank $r$, then
\[
S(\CCPi{h}) = \log r.
\]
\end{lemma}

\begin{proof}
The spectrum of $\rho_h = \CCPi{h} / r$ is $\{1/r\}$ with multiplicity $r$ and $\{0\}$ otherwise.
Thus $-\Tr(\rho_h \log \rho_h) = -r \cdot (1/r) \log(1/r) = \log r$.
\end{proof}

\begin{remark}[Entropy and horizon refinement]
\lean{CoherenceCalculus.AppendixOperators.entropy_horizon_refinement}
For nested horizons $h \le h'$ with $\rank(\CCPi{h}) \le \rank(\CCPi{h'})$, we have $S(\CCPi{h}) \le S(\CCPi{h'})$.
This is consistent with the telescoping identity (HFT-2, Theorem~\ref{thm:HFT}): finer horizons have greater representational capacity.
\end{remark}

\begin{definition}[Defect-corrected entropy (optional refinement)]
\label{def:defect-entropy}\lean{CoherenceCalculus.AppendixOperators.defectCorrectedEntropy}
For a state $x \in \cH$ with $\CCPi{h} x \ne 0$, and defect energy operator $\CCDefOp{d}{h} = \CCLeak{d}{h}^\ast \CCLeak{d}{h}$, define the \emph{defect-corrected entropy}
\[
S_{\mathrm{def},h}(x) := S(\CCPi{h}) + \lambda \frac{\langle x, \CCPi{h} \CCDefOp{d}{h} \CCPi{h} x \rangle}{\langle x, \CCPi{h} x \rangle},
\]
where $\lambda \ge 0$ is a weighting parameter.
This adds a penalty proportional to the normalized defect cost of $x$.
\end{definition}

\begin{quote}\small
\textbf{Intuition.}
Horizon entropy measures the log-count of observable degrees of freedom ($\log r$).
The optional defect correction penalizes states that ``leak'' into the shadow, providing a combined capacity-plus-cost functional.
\end{quote}

% ------------------------------------------------------------
\subsection{Topological Spectrograph (Graded Invariants from Complexes)}
\label{subsec:topological-spectrograph}
% ------------------------------------------------------------

This subsection records additional structure built on cohomology.
For the base definitions of Hilbert complexes, cohomology $\CCCoh{k}(D)$, Hodge Laplacian, Green operator, and homotopy operator, see \S\ref{subsec:solver-layer}.

\begin{definition}[Topological spectrograph]
\label{def:spectrograph}\lean{CoherenceCalculus.AppendixOperators.topologicalSpectrograph}
The \emph{topological spectrograph} of a complex $(C^\bullet, D)$ is the sequence of cohomology dimensions:
\[
\mathrm{SpecTop}(C^\bullet) := \bigl( \dim \CCCoh{k}(D) \bigr)_{k}.
\]
\end{definition}

\begin{definition}[Graded volume functional]
\label{def:graded-volume}\lean{CoherenceCalculus.AppendixOperators.gradedVolume}
For a weight sequence $w = (w_k)_{k \ge 0}$ with $w_k \ge 0$, define the \emph{weighted graded volume}
\[
V_w(C^\bullet) := \sum_k w_k \dim \CCCoh{k}(D).
\]
Natural choices: $w_k = 1$ (total dimension) or $w_k = (-1)^k$ (Euler characteristic $\chi$).
\end{definition}

\begin{lemma}[Additivity under direct sums]
\label{lem:spectrograph-additivity}\lean{CoherenceCalculus.AppendixOperators.spectrograph_additivity}
For complexes $(C^\bullet, D)$ and $(C'^\bullet, D')$:
\begin{enumerate}[label=(\roman*)]
\item $\mathrm{SpecTop}(C^\bullet \oplus C'^\bullet) = \mathrm{SpecTop}(C^\bullet) + \mathrm{SpecTop}(C'^\bullet)$ (componentwise).
\item $V_w(C^\bullet \oplus C'^\bullet) = V_w(C^\bullet) + V_w(C'^\bullet)$.
\end{enumerate}
\end{lemma}

\begin{proof}
Cohomology respects direct sums: $\CCCoh{k}(C \oplus C') \cong \CCCoh{k}(C) \oplus \CCCoh{k}(C')$.
Dimensions add, giving (i). Linearity of summation gives (ii).
\end{proof}

\begin{remark}[Euler characteristic and exactness]
\lean{CoherenceCalculus.AppendixOperators.exact_complex_spectrograph_zero}
The Euler characteristic $\chi(C^\bullet) := \sum_k (-1)^k \dim \CCCoh{k}(D)$ vanishes if and only if the alternating sum of Betti numbers is zero.
For an exact complex ($\CCCoh{k}(D) = 0$ for all $k$), trivially $\chi = 0$ and the spectrograph is the zero sequence.
\end{remark}

% ------------------------------------------------------------
\subsection{Flow--Stability Indices (Fixed Points of Iteration Maps)}
\label{subsec:flow-stability}
% ------------------------------------------------------------

Stability analysis concerns fixed points of iteration maps, without reference to continuous parameter evolution.

\begin{definition}[Iteration map and fixed point]
\label{def:iteration-map}\lean{CoherenceCalculus.AppendixOperators.IterationMap}\lean{CoherenceCalculus.AppendixOperators.isFixedPoint}
Let $F_h: X_h \to X_h$ be a map on an observable state space $X_h$ (finite-dimensional for accessibility).
A point $x^* \in X_h$ is a \emph{fixed point} if $F_h(x^*) = x^*$.
\end{definition}

\begin{definition}[Linearization at a fixed point]
\label{def:linearization}\lean{CoherenceCalculus.AppendixOperators.Linearization}
If $F_h$ is differentiable at a fixed point $x^*$, the \emph{linearization} is the Jacobian
\[
J := DF_h(x^*).
\]
For a linear map $F_h$, the linearization is $F_h$ itself.
\end{definition}

\begin{definition}[Stability index]
\label{def:stability-index}\lean{CoherenceCalculus.AppendixOperators.stabilityIndex}
The \emph{stability index} of a fixed point $x^*$ is
\[
\mathrm{Ind}_{\mathrm{stab}}(x^*) := \#\{ \text{eigenvalues } \lambda \text{ of } J \text{ with } |\lambda| > 1 \},
\]
counted with algebraic multiplicity.
\end{definition}

\begin{definition}[Linearized stable and unstable types]
\label{def:stable-unstable}\lean{CoherenceCalculus.AppendixOperators.isStableFixedPoint}\lean{CoherenceCalculus.AppendixOperators.isUnstableFixedPoint}
A fixed point $x^*$ is of:
\begin{itemize}[itemsep=0.2em]
\item \emph{linearly stable type} if the spectral radius $\rho(J) < 1$, i.e., all eigenvalues satisfy $|\lambda| < 1$.
\item \emph{linearly unstable type} if $\rho(J) > 1$, i.e., at least one eigenvalue satisfies $|\lambda| > 1$.
\end{itemize}
This is a classification of the Jacobian \(J = DF_h(x^*)\). Turning it into a
nonlinear local-attractor theorem requires additional hypotheses and standard
results from discrete dynamical systems, which are not developed here.
\end{definition}

\begin{lemma}[Exact criterion for linear iteration maps]
\label{lem:stability-criterion}\lean{CoherenceCalculus.AppendixOperators.stability_criterion}\lean{CoherenceCalculus.AppendixOperators.instability_criterion}
Assume \(X_h\) is finite-dimensional and \(F_h : X_h \to X_h\) is linear.
Let \(x^* = 0\), so \(J = F_h\).
Then:
\begin{enumerate}[label=(\roman*)]
\item if \(\rho(J) < 1\), then \(0\) is globally attracting and
      \(F_h^n(x_0) \to 0\) for every \(x_0 \in X_h\);
\item if \(\rho(J) > 1\), then \(0\) is unstable, and every initial condition
      with nonzero component in the direct sum of generalized eigenspaces for
      eigenvalues \(|\lambda| > 1\) has iterates whose norms do not converge to
      \(0\).
\end{enumerate}
\end{lemma}

\begin{proof}
Because \(X_h\) is finite-dimensional, \(J\) admits a Jordan normal form over
\(\mathbb C\). Every Jordan block for an eigenvalue \(\lambda\) contributes a
factor of the form \(n^m \lambda^n\) to \(J^n\).

If \(\rho(J) < 1\), then every eigenvalue satisfies \(|\lambda| < 1\), so
every such factor tends to \(0\). Hence \(J^n \to 0\), and therefore
\(F_h^n(x_0) = J^n x_0 \to 0\) for every \(x_0\).

If \(\rho(J) > 1\), then there is at least one generalized eigenspace
corresponding to an eigenvalue with \(|\lambda| > 1\). On any nonzero vector in
that generalized eigenspace, the Jordan-block formula shows that the iterates
cannot converge to \(0\). Therefore \(0\) is unstable.
\end{proof}

\begin{remark}[Connection to admissibility]
\lean{CoherenceCalculus.AppendixOperators.stabilityIndex}
In the context of admissibility and branch selection (Section~\ref{sec:admissibility}), the selected branch corresponds to the stable fixed point of an associated iteration.
The stability index counts ``unstable directions,'' namely the number of
constraints violated by perturbations.
\end{remark}

\begin{remark}[Nonlinear extension]
\lean{CoherenceCalculus.AppendixOperators.instability_criterion}
For genuinely nonlinear iteration maps, promoting the linearized type of
Definition~\ref{def:stable-unstable} to a local-attractor or local-instability
statement requires additional hypotheses and standard discrete-dynamical
results such as the stable-manifold or Hartman--Grobman theorems. Those
extensions are outside the present appendix.
\end{remark}

\begin{example}[Linear iteration in $\bbR$]
\lean{CoherenceCalculus.AppendixOperators.linear_stability_example}
\label{ex:linear-stability}
Let $F(x) = ax$ for $a \in \bbR$.
The fixed point is $x^* = 0$ with $J = a$.
\begin{itemize}[itemsep=0.2em]
\item $|a| < 1$: stable, $\mathrm{Ind}_{\mathrm{stab}} = 0$.
\item $|a| > 1$: unstable, $\mathrm{Ind}_{\mathrm{stab}} = 1$.
\end{itemize}
\end{example}

\begin{quote}\small
\textbf{Intuition.}
Stability here is about iteration under an admissibility or solver map, not about evolution in an external parameter.
The stability index counts the number of ``expanding directions'' at a fixed point.
\end{quote}

\begin{remark}[Alternative indices]
\lean{CoherenceCalculus.AppendixOperators.alternative_stability_indices}
For more refined invariants (e.g., Morse index, Conley index), additional structure is required.
These are not developed here but follow similar principles: counting dimensions of unstable/stable manifolds.
\end{remark}
